\chapter{การตรวจวัดเชิงกายภาพและการบินเชิงปริมาณ}
\label{chap:ch04}

% --- หน้าที่ 1 ของบท: แผนบริหารการสอน (กล่อง 1 และ 2) ---
\section*{แผนบริหารการสอนประจำบทที่ 4}
\addcontentsline{toc}{section}{แผนบริหารการสอนประจำบทที่ 4}

\begin{planbox}{หัวข้อเนื้อหาประจำบท}
\begin{enumerate}
    \item การคำนวณระยะตัวอย่างภาคพื้นดิน GSD และการคำนวณขนาดพื้นที่จริง
    \item การแปลงหน่วยพื้นที่ตามมาตราวัดไทย ไร่ งาน และตารางวา
    \item อัลกอริทึมการนับจำนวนต้นพืชและระบุพิกัดกึ่งกลางทรงพุ่ม
    \item การตรวจจับช่องว่างและจุดเว้าแหว่งของแปลงเพื่อการวางแผนปลูกซ่อม
    \item การคำนวณอัตราการพ่นสารชีวภัณฑ์และการไหลของของเหลว
    \item การวิเคราะห์ประสิทธิภาพการใช้พลังงานและเวลาบินคงเหลือ
\end{enumerate}
\end{planbox}

\begin{planbox}{วัตถุประสงค์เชิงพฤติกรรม}
\begin{enumerate}
    \item คำนวณค่าความละเอียดภาคพื้นดิน GSD จากความสูงบินและคุณลักษณะเลนส์ได้อย่างถูกต้อง
    \item แปลงค่าพื้นที่จากตารางเมตรเป็นหน่วยไร่ งาน ตารางวาได้อย่างแม่นยำ
    \item อธิบายอัลกอริทึมตรวจนับต้นพืชด้วยจุดยอดสัมพัทธ์ Local Maxima ได้อย่างชัดเจน
    \item ประเมินจำนวนต้นกล้าที่ต้องปลูกซ่อมจากอัตราจุดเว้าแหว่งในแปลงได้อย่างมีประสิทธิภาพ
    \item คำนวณอัตราการไหลของหัวฉีดและปริมาตรสารชีวภัณฑ์ต่อไร่ได้อย่างสมบูรณ์
\end{enumerate}
\end{planbox}

\newpage

% --- หน้าที่ 2 ของบท: แผนบริหารการสอน (กล่อง 3, 4, 5) ---
\begin{planbox}{วิธีสอนและกิจกรรมการเรียนการสอน}
\begin{enumerate}
    \item การบรรยายหลักการทางเรขาคณิตของการถ่ายภาพทางอากาศและคณิตศาสตร์การเกษตร
    \item การทดลองจำลองการนับต้นไม้ในแปลงทุเรียนด้วยโปรแกรมจำลองเสมือน
    \item การฝึกคำนวณอัตราการพ่นสารชีวภัณฑ์เทียบกับความเร็วการบินภาคพื้น
    \item การทำโครงงานย่อยประเมินต้นทุนพลังงานและประสิทธิภาพการบินต่อพื้นที่หนึ่งไร่
\end{enumerate}
\end{planbox}

\begin{planbox}{สื่อการเรียนการสอน}
\begin{enumerate}
    \item โปรแกรมจำลองการคำนวณขนาดพื้นที่และอัตราการพ่นสารในระบบ Flutter
    \item ชุดภาพถ่ายทางอากาศแปลงไม้ผลที่มีข้อมูลพิกัดทางภูมิศาสตร์อ้างอิง
    \item แผนภาพแสดงเรขาคณิตของกรวยภาพกล้องและระยะตกกระทบภาคพื้นดิน
    \item เครื่องคิดเลขวิทยาศาสตร์และเอกสารสูตรการคำนวณเชิงปริมาณ
\end{enumerate}
\end{planbox}

\begin{planbox}{การวัดและประเมินผล}
\begin{enumerate}
    \item การประเมินผลการทำโจทย์คำนวณทางคณิตศาสตร์การบินและมาตราวัดไทย
    \item การประเมินความแม่นยำของผลการนับจำนวนต้นพืชเมื่อเทียบกับข้อมูลภาคพื้นดิน
    \item การส่งรายงานการประเมินการใช้สารเคมีและพลังงานแบตเตอรี่ประจำบท
\end{enumerate}
\end{planbox}

\clearpage

% --- หน้าที่ 3 เป็นต้นไป: เนื้อหาวิชาการ ---
\section{การคำนวณระยะตัวอย่างภาคพื้นดินและขนาดพื้นที่จริง}

การแปลงข้อมูลจากพิกัดพิกเซลบนภาพถ่ายดิจิทัลไปสู่มิติทางกายภาพจริงบนพื้นผิวโลก จำเป็นต้องคำนวณค่า \textbf{ระยะตัวอย่างภาคพื้นดิน (Ground Sample Distance หรือ GSD)} ซึ่งแสดงความยาวจริงบนพื้นดินที่แทนด้วยขนาดของจุดภาพหนึ่งพิกเซล (มีหน่วยเป็นเซนติเมตรต่อพิกเซล)

\begin{formulabox}[สมการคำนวณระยะตัวอย่างภาคพื้นดิน GSD]
\begin{equation}
\text{GSD} = \frac{H \cdot S_w}{f \cdot I_w} \times 100
\label{eq:gsd}
\end{equation}
โดยที่
$H$ คือความสูงของอากาศยานเหนือผิวดิน (เมตร)\newline
$S_w$ คือความกว้างของเซนเซอร์รับภาพ (มิลลิเมตร)\newline
$f$ คือความยาวโฟกัสของเลนส์กล้อง (มิลลิเมตร)\newline
$I_w$ คือความกว้างของภาพถ่ายดิจิทัล (พิกเซล)
\end{formulabox}

เมื่อทราบค่า GSD และขนาดความกว้างความยาวของภาพถ่าย ($I_w \times I_h$) จะสามารถคำนวณขนาดพื้นที่ครอบคลุมของภาพถ่ายบนพื้นผิวโลกจริงได้ตามสมการ

\begin{equation}
\text{Area} = \left(\frac{I_w \cdot \text{GSD}}{100}\right) \times \left(\frac{I_h \cdot \text{GSD}}{100}\right) \quad (\text{ตารางเมตร})
\end{equation}

\section{การแปลงหน่วยพื้นที่ตามมาตราวัดไทย}

ในบริบทของการเกษตรกรรมในประเทศไทย การสื่อสารขนาดแปลงที่ดินนิยมใช้หน่วยวัดประเพณีไทย ได้แก่ \textbf{ไร่ งาน และตารางวา} มากกว่าหน่วยเอเคอร์หรือเฮกตาร์ ระบบจึงแปลงค่าพื้นที่จากตารางเมตรตามเกณฑ์มาตรฐานทางกฎหมายที่ดินของไทย

\begin{table}[H]
\caption{อัตราส่วนการแปลงหน่วยพื้นที่ตามมาตราวัดที่ดินของไทย}
\label{tab:thai_land_units}
\begin{tabularx}{\textwidth}{p{3.2cm}p{3.2cm}p{3.2cm}Y}
\toprule
\textbf{ชื่อหน่วยวัด} & \textbf{ขนาดในหน่วยตารางเมตร} & \textbf{สัดส่วนเทียบหน่วยไร่} & \textbf{การนำไปใช้ในระบบ} \\
\midrule
1 ตารางวา (ตร.ว.) & 4 ตารางเมตร & $1/400$ ไร่ & หน่วยย่อยเศษทศนิยม \\
1 งาน & 400 ตารางเมตร & $1/4$ ไร่ & สัดส่วนย่อยขนาดกลาง \\
1 ไร่ & 1,600 ตารางเมตร & 1.0 ไร่ & หน่วยหลักในการเกษตรไทย \\
1 เฮกตาร์ (ha) & 10,000 ตารางเมตร & 6.25 ไร่ & มาตรฐานสากลสำหรับวิจัย \\
\bottomrule
\end{tabularx}
\end{table}

\section{อัลกอริทึมการนับจำนวนต้นพืชและระบุพิกัดกึ่งกลาง}

การนับจำนวนต้นไม้ในแปลงผลไม้เศรษฐกิจ เช่น ทุเรียน มังคุด หรือสวนปาล์มน้ำมัน อาศัยอัลกอริทึมการตรวจหาจุดยอดสัมพัทธ์ในกริดระนาบ \textbf{Local Maxima Centroid Detection} โดยมีขั้นตอนการประมวลผลดังนี้
\begin{enumerate}
    \item สแกนค้นหาจุดภาพที่มีค่าดัชนี VARI หรือ ExG สูงกว่าเกณฑ์ความไวทรงพุ่ม
    \item ค้นหาจุดที่มีค่าสูงสุดในขอบเขตบริเวณใกล้เคียง (Search Window) ขนาด $16 \times 16$ พิกเซล เพื่อหาจุดยอดพุ่มไม้
    \item บันทึกพิกัดกึ่งกลาง $(X_c, Y_c)$ ของต้นไม้ พร้อมตรวจจับระยะห่างเพื่อไม่ให้เกิดการนับซ้ำซ้อนในต้นเดียวกัน
    \item นำพิกัดไปวาดวงกลมสัญลักษณ์มาร์กเกอร์สีเหลืองทับบนภาพถ่ายเพื่อให้นักบินสามารถตรวจสอบด้วยสายตาได้ทันที
\end{enumerate}

\section{การตรวจจับช่องว่างของแปลงเพื่อการวางแผนปลูกซ่อม}

ในแปลงเกษตรกรรมแบบร่องสวน ต้นไม้มักจะถูกปลูกด้วยระยะห่างที่สม่ำเสมอ เช่น $8 \times 8$ เมตรในสวนทุเรียน อัลกอริทึมจะตรวจสอบโครงสร้างกริดของแนวปลูก หากพบว่าตำแหน่งใดในกริดไม่มีจุดยอดของทรงพุ่มพืช (ค่าดัชนีสเปกตรัมต่ำกว่าเกณฑ์) ระบบจะบันทึกเป็น \textbf{จุดเว้าแหว่งหรือต้นที่ขาดหาย (Missing Plant Gap)}

ข้อมูลจุดเว้าแหว่งนี้จะถูกสรุปเป็นจำนวนต้นกล้าที่แนะนำให้เกษตรกรจัดเตรียมเพื่อนำไปปลูกซ่อมแซมได้อย่างแม่นยำ โดยไม่ต้องเดินสำรวจด้วยเท้าทั่วทั้งแปลง

\section{การคำนวณอัตราการพ่นสารชีวภัณฑ์และการไหลของของเหลว}

สำหรับโดรนการเกษตรที่มีการติดตั้งถังบรรจุของเหลวและหัวฉีดพ่น การควบคุมอัตราการพ่นให้สม่ำเสมอต่อหน่วยพื้นที่ ถือเป็นปัจจัยชี้วัดคุณภาพของการให้ปุ๋ยชีวภาพและสารควบคุมศัตรูพืช อัตราการไหลของหัวฉีด (Flow Rate) มีความสัมพันธ์กับความเร็วการบินของโดรนและความกว้างของแถบพ่น

\begin{formulabox}[สมการคำนวณอัตราการไหลของหัวฉีดพ่น]
\begin{equation}
Q = \frac{R \cdot v \cdot W \cdot 60}{1,600}
\label{eq:flow_rate}
\end{equation}
โดยที่
$Q$ คืออัตราการไหลรวมของหัวฉีด (ลิตรต่อนาที)\newline
$R$ คืออัตราการพ่นสารเป้าหมายที่ต้องการ (ลิตรต่อไร่)\newline
$v$ คือความเร็วการบินภาคพื้น (เมตรต่อวินาที)\newline
$W$ คือความกว้างของแถบพ่นสาร (เมตร)\newline
1,600 คือจำนวนตารางเมตรต่อหนึ่งไร่
\end{formulabox}

หากโดรนเร่งความเร็วการบินขึ้น ระบบควบคุมสมองกลจะเพิ่มสัญญาณ PWM ไปยังปั๊มพ่นโดยอัตโนมัติ เพื่อรักษาอัตราการพ่นต่อไร่ให้คงที่สม่ำเสมอทั่วทั้งแปลง

\section{การวิเคราะห์ประสิทธิภาพการใช้พลังงานและเวลาบินคงเหลือ}

ระบบควบคุมจะตรวจวัดอัตราการสิ้นเปลืองพลังงานไฟฟ้าจากแบตเตอรี่ลิเทียมพอลิเมอร์ โดยคำนวณเป็นอัตราส่วนวัตต์-ชั่วโมงต่อไร่ (Watt-hours per Rai) ซึ่งอากาศยานไร้คนขับสำหรับการเกษตรทั่วไปจะมีอัตราการใช้พลังงานเฉลี่ยอยู่ที่ประมาณ 60 ถึง 75 วัตต์-ชั่วโมงต่อไร่

จากนั้นระบบจะประเมินเวลาการบินคงเหลือ (Estimated Remaining Flight Time) เทียบกับเปอร์เซ็นต์ประจุไฟฟ้าที่เหลืออยู่ โดยหากเวลาบินคงเหลือเริ่มใกล้เคียงกับเวลาที่ต้องใช้ในการบินกลับสู่จุดปล่อยตัว (Return to Home หรือ RTH) ระบบจะส่งสัญญาณเตือนให้นักบินนำอากาศยานกลับฐานเพื่อเปลี่ยนแบตเตอรี่อย่างปลอดภัย

\clearpage

\section*{คำถามทบทวนท้ายบทที่ 4}
\addcontentsline{toc}{section}{คำถามทบทวนท้ายบทที่ 4}

\begin{enumerate}
    \item โดรนติดตั้งกล้องความยาวโฟกัส 4.0 มิลลิเมตร ขนาดเซนเซอร์กว้าง 6.0 มิลลิเมตร บินที่ระดับความสูง 20 เมตรเหนือแปลงเกษตร หากภาพถ่ายมีความกว้าง 4,000 พิกเซล จงคำนวณค่า GSD ในหน่วยเซนติเมตรต่อพิกเซล
    \item ภาพถ่ายทางอากาศแปลงหนึ่งครอบคลุมพื้นที่ 8,000 ตารางเมตร จงแปลงค่าเป็นหน่วย ไร่ งาน และตารางวา ตามมาตราวัดไทย
    \item เกษตรกรต้องการพ่นปุ๋ยอินทรีย์ทางใบในอัตรา 5 ลิตรต่อไร่ โดยใช้โดรนบินด้วยความเร็ว 4 เมตรต่อวินาที และมีความกว้างแถบพ่น 3 เมตร จงคำนวณอัตราการไหลของปั๊มที่ต้องตั้งค่าในหน่วยลิตรต่อนาที
    \item อัลกอริทึม Local Maxima Centroid Detection มีข้อดีอย่างไรในการนับต้นไม้ในสวนผลไม้เมื่อเทียบกับการคำนวณเพียงพื้นที่สีเขียวรวม
\end{enumerate}

\clearpage
